\documentclass[../../../include/open-logic-section]{subfiles}

\begin{document}

\olfileid{sth}{ordinals}{basic}
\olsection{ویژگی‌های بنیادی اعداد ترتیبی}

دیدیم که نخستین چند عدد ترتیبی همان اعداد طبیعی‌اند. دلیل اصلی پروراندن
نظریه‌ای دربارهٔ اعداد ترتیبی، گسترش اصل استقرایی است که دربارهٔ اعداد
طبیعی برقرار است. با دنباله‌ای از نتایج مقدماتی به‌تدریج به این هدف
خواهیم رسید.

\begin{lem}\ollabel{ordmemberord}
هر~!!{element} یک عدد ترتیبی، خود عدد ترتیبی است.
\end{lem}

\begin{proof}
فرض کنید~$\alpha$ عدد ترتیبی و~$b \in \alpha$ باشد. چون~$\alpha$ متعدی
است، $b \subseteq \alpha$. پس همان‌گونه که~$\in$، $\alpha$ را خوش‌ترتیب
می‌کند، $b$ را نیز خوش‌ترتیب می‌کند.

برای دیدن اینکه~$b$ متعدی است، فرض کنید~$x \in c \in b$. چون
$b \subseteq \alpha$، داریم~$c \in \alpha$. باز چون~$\alpha$ متعدی
است، $c \subseteq \alpha$ و در نتیجه~$x \in \alpha$. پس
$x,c,b \in \alpha$. اما~$\in$، $\alpha$ را خوش‌ترتیب می‌کند؛ بنابراین
بنا بر~\olref[sth][ordinals][wo]{wo:strictorder}، $\in$ رابطه‌ای متعدی
روی~$\alpha$ است. پس از~$x \in c \in b$ نتیجه می‌شود~$x \in b$. با
تعمیم، $c \subseteq b$.
\end{proof}

\begin{cor}\ollabel{ordissetofsmallerord}
برای هر عدد ترتیبی~$\alpha$،
$\alpha = \Setabs{\beta \in \alpha}{\beta \text{ یک عدد ترتیبی است}}$.
\end{cor}

\begin{proof}
بی‌درنگ از~\olref{ordmemberord} نتیجه می‌شود.
\end{proof}

مضمون تقریبی دو نتیجهٔ اصلی بعدی، یعنی~\olref{ordinductionschema} و
\olref{ordtrichotomy}، این است که خود اعداد ترتیبی با عضویت خوش‌ترتیب
شده‌اند:

\begin{thm}[استقرای ترامتناهی]\ollabel{ordinductionschema}
برای هر فرمول~$\phi(x)$:
\[
	\text{اگر }\exists \alpha \phi(\alpha)\text{، آنگاه }\exists \alpha(\phi(\alpha)
	\land (\forall \beta \in \alpha) \lnot \phi(\beta))
\]
که در آن، سورهای نمایش‌داده‌شده به‌طور ضمنی به اعداد ترتیبی محدودند.
\end{thm}

\begin{proof}
%\olref{ordinductionschema1}.
فرض کنید~$\phi(\alpha)$ برای یک عدد ترتیبی~$\alpha$. اگر
$(\forall \beta \in \alpha)\lnot\phi(\beta)$، کار تمام است. در غیر این
صورت، مجموعهٔ
$X=\Setabs{\beta\in\alpha}{\phi(\beta)}$ ناتهی است. چون~$\in$،
$\alpha$ را خوش‌ترتیب می‌کند، $X$ یک~!!{element}~$\in$-کمینه مانند
$\gamma$ دارد. بنا بر~\olref{ordmemberord}، $\gamma$ عدد ترتیبی است؛ و
چون~$\gamma\in X$، داریم~$\phi(\gamma)$. همچنین کمینگی~$\gamma$ و تعدی
$\alpha$ نتیجه می‌دهند
$(\forall\beta\in\gamma)\lnot\phi(\beta)$. پس~$\gamma$ همان عدد ترتیبی
مطلوب است.
%	\olref{ordinductionschema2}. فرض کنید یک عدد ترتیبی~$\gamma$ وجود دارد
%	به‌طوری‌که $\lnot\phi(\gamma)$. آنگاه بنا بر
%	\olref{ordinductionschema1}، یک عدد ترتیبی~$\in$-مینیمال~$\alpha$
%	وجود دارد که $\lnot\phi(\alpha)$. پس $(\forall \beta <
%	\alpha)\phi(\beta)$ و در نتیجه مقدم شرطی کاذب است.
\end{proof}
\noindent
توجه کنید که~\olref{ordinductionschema} را به‌طور هم‌ارز می‌توان به صورت
طرحوارهٔ زیر بیان کنیم؛ کافی است~$\lnot\phi(\alpha)$ را
در~\olref{ordinductionschema} قرار دهیم و سپس دستکاری‌های مقدماتی منطقی انجام دهیم:
\[
\text{اگر }\forall \alpha((\forall \beta \in \alpha)\phi(\beta) \lif
\phi(\alpha))\text{، آنگاه }\forall \alpha\phi(\alpha).
\]

\begin{thm}[سه‌شقی]\ollabel{ordtrichotomy}
برای هر دو عدد ترتیبی~$\alpha$ و~$\beta$،
$\alpha \in \beta \lor \alpha = \beta \lor \beta \in \alpha$.
\end{thm}

\begin{proof}
اثبات با استقرای دوگانه است؛ یعنی دو بار از~\olref{ordinductionschema}
استفاده می‌کنیم. می‌گوییم~$x$ با~$y$ \emph{قابل مقایسه} است اگر و تنها
اگر~$x \in y \lor x=y \lor y \in x$.

برای استقرا، فرض کنید هر عدد ترتیبیِ عضو~$\alpha$ با \emph{هر} عدد
ترتیبی قابل مقایسه باشد. برای استقرای دوم، فرض کنید~$\alpha$ با هر عدد
ترتیبیِ عضو~$\beta$ قابل مقایسه باشد. نشان خواهیم داد~$\alpha$ با~$\beta$
قابل مقایسه است. با استقرا روی~$\beta$ نتیجه خواهد شد که~$\alpha$ با هر
عدد ترتیبی قابل مقایسه است؛ و سپس با استقرا روی~$\alpha$، \emph{هر} عدد
ترتیبی با \emph{هر} عدد ترتیبی قابل مقایسه خواهد بود، چنان‌که مطلوب است.
کافی است فرض کنیم~$\alpha \notin \beta$ و~$\beta \notin \alpha$ و نشان
دهیم~$\alpha=\beta$.

برای نشان‌دادن~$\alpha \subseteq \beta$، یک~$\gamma \in \alpha$ ثابت
بگیرید؛ بنا بر~\olref{ordmemberord}، $\gamma$ عدد ترتیبی است. پس بنا بر
فرض استقرای نخست، $\gamma$ با~$\beta$ قابل مقایسه است. اما اگر
$\gamma=\beta$ یا~$\beta \in \gamma$، آنگاه~$\beta \in \alpha$ (در صورت
لزوم با استفاده از تعدی~$\alpha$)، برخلاف فرض ما. پس~$\gamma \in \beta$.
با تعمیم، $\alpha \subseteq \beta$.

استدلالی کاملاً مشابه با استفاده از فرض استقرای دوم نشان می‌دهد
$\beta \subseteq \alpha$. پس~$\alpha=\beta$.
\end{proof}\noindent
ازاین‌رو گاه به‌جای~$\alpha \in \beta$ می‌نویسیم~$\alpha<\beta$، زیرا
$\in$ مانند یک رابطهٔ ترتیب رفتار می‌کند. جز آشنایی و این واقعیت که نوشتن
$\alpha \leq \beta$ از~$\alpha \in \beta \lor \alpha=\beta$ آسان‌تر
است، دلیل ژرفی برای این کار وجود ندارد.\footnote{می‌توانستیم بنویسیم
$\alpha \mathrel{\underline{\in}} \beta$، اما این نمادگذاری کاملاً
نامتعارف بود.}

در ادامه دو پیامد سریع از نتایج اخیر می‌آید که نخستین آن‌ها نمادگذاری تازه
را به کار می‌گیرد:

\begin{cor}\ollabel{ordordered}
اگر~$\exists \alpha\phi(\alpha)$، آنگاه
$\exists \alpha(\phi(\alpha) \land \forall \beta(\phi(\beta) \lif
\alpha \leq \beta))$. افزون بر این، برای هر اعداد ترتیبی
$\alpha,\beta,\gamma$، هم~$\alpha \notin \alpha$ و هم
$\alpha \in \beta \in \gamma \lif \alpha \in \gamma$.
\end{cor}

\begin{proof}
درست مانند~\olref[wo]{wo:strictorder}.
\end{proof}

\begin{prob}
«استدلال کاملاً مشابه» در اثبات
\olref[sth][ordinals][basic]{ordtrichotomy} را کامل کنید.
\end{prob}

\begin{cor}\ollabel{corordtransitiveord}
$A$ عدد ترتیبی است اگر و تنها اگر مجموعه‌ای متعدی از اعداد ترتیبی باشد.
\end{cor}

\begin{proof}
\emph{از چپ به راست.} بنا بر~\olref{ordmemberord}.
\emph{از راست به چپ.} اگر~$A$ مجموعه‌ای متعدی از اعداد ترتیبی باشد، آنگاه
بنا بر~\olref{ordinductionschema} و~\olref{ordtrichotomy}، $\in$ مجموعهٔ
$A$ را خوش‌ترتیب می‌کند.
\end{proof}

اکنون~\olref{ordinductionschema} و~\olref{ordtrichotomy} را چنین تفسیر
کردیم که~$\in$ اعداد ترتیبی را خوش‌ترتیب می‌کند. اما بنا بر نتیجهٔ زیر،
باید دربارهٔ چنین ادعایی \emph{بسیار محتاط} باشیم:

\begin{thm}[پارادوکس بورالی--فورتی]\ollabel{buraliforti}
هیچ مجموعه‌ای از همهٔ اعداد ترتیبی وجود ندارد.
\end{thm}

\begin{proof}
برای برهان خلف، فرض کنید~$O$ مجموعهٔ همهٔ اعداد ترتیبی باشد. اگر
$\alpha \in \beta \in O$، آنگاه بنا بر~\olref{ordmemberord}، $\alpha$
عدد ترتیبی است و بنابراین~$\alpha \in O$. پس~$O$ متعدی است و در نتیجه
بنا بر~\olref{corordtransitiveord}، $O$ یک عدد ترتیبی است. ازاین‌رو
$O \in O$، برخلاف~\olref{ordordered}.
\end{proof}
\noindent
این نتیجه به نام~\citeauthor{Burali-Forti1897} نام‌گذاری شده است. اما
کانتور بود که در سال ۱۸۹۹---در نامه‌ای به ددکیند---نخستین بار
\emph{تناقض} موجود در فرض مجموعه‌ای از همهٔ اعداد ترتیبی را به‌روشنی دید.
چنان‌که فان هاینورت توضیح می‌دهد:
\begin{quote}
  خود بورالی--فورتی تناقض را مؤید این نتیجه می‌دانست که ترتیب طبیعی اعداد
  ترتیبی صرفاً ترتیبی جزئی است، و این نتیجه را با \emph{برهان خلف} به دست
  می‌آورد.
  \citep[p.~105]{Heijenoort1967}
\end{quote}
اگر اشتباه بورالی--فورتی را کنار بگذاریم، می‌توانیم مطالب پیشین را چنین
خلاصه کنیم: اعداد ترتیبی مجموعه‌هایی‌اند که به‌طور منفرد با عضویت
خوش‌ترتیب شده‌اند و در مجموع نیز با عضویت خوش‌ترتیب‌اند (بی‌آنکه در مجموع
یک مجموعه را تشکیل دهند).

برای تکمیل بحث، چند ویژگی بنیادی دیگر اعداد ترتیبی را می‌آوریم که از
\olref{ordinductionschema} و~\olref{ordtrichotomy} نتیجه می‌شوند.

\begin{prop}
هر دنبالهٔ اکیداً نزولی از اعداد ترتیبی متناهی است.
\end{prop}

\begin{proof}
هر دنبالهٔ نامتناهی و اکیداً نزولی از اعداد ترتیبی
$\alpha_0>\alpha_1>\alpha_2>\ldots$ هیچ عضو~$<$-مینیمالی ندارد، برخلاف
\olref{ordinductionschema}.
\end{proof}

\begin{prop}\ollabel{ordinalsaresubsets}
برای هر دو عدد ترتیبی~$\alpha,\beta$،
$\alpha \subseteq \beta \lor \beta \subseteq \alpha$.
\end{prop}

\begin{proof}
اگر~$\alpha \in \beta$، چون~$\beta$ متعدی است،
$\alpha \subseteq \beta$. به همین ترتیب، اگر~$\beta \in \alpha$، آنگاه
$\beta \subseteq \alpha$. و اگر~$\alpha=\beta$، آنگاه هم
$\alpha \subseteq \beta$ و هم~$\beta \subseteq \alpha$. پس با
\olref{ordtrichotomy} کار تمام است.
\end{proof}

\begin{prop}\ollabel{ordisoidentity}
برای هر دو عدد ترتیبی~$\alpha,\beta$، $\alpha=\beta$ اگر و تنها اگر
$\ordeq{\alpha}{\beta}$.
\end{prop}

\begin{proof}
اعداد ترتیبی خوش‌ترتیب‌اند؛ پس این نتیجه بی‌درنگ از سه‌شقی
(\olref[basic]{ordtrichotomy}) و~\olref[iso]{wellordnotinitial} به دست
می‌آید.
\end{proof}

\begin{prob}\ollabel{probunionordinalsordinal}
اثبات کنید اگر هر عضو~$X$ یک عدد ترتیبی باشد، آنگاه~$\bigcup X$ یک عدد
ترتیبی است.
\end{prob}

\end{document}
